\section{Results and Evidence Required for Claims}
\label{sec:results}

\subsection{Current status}

There are no publishable RL results in the supplied codebase. The existing manuscript tables that report 2018--2026 benchmark performance are not generated by the current corrected RL protocol, omit the proposed model, and should not appear in a submission. Likewise, the ablation and sensitivity scripts currently manufacture ``plausible'' metrics through hard-coded degradation factors and random noise. Such outputs are not simulations of the implemented model and must never be reported as empirical evidence.

This section is intentionally a results-ready reporting structure. It becomes a results section only after the corrections in \cref{sec:methodology} and the evaluation programme in \cref{sec:data} are completed.

\subsection{Required primary table}

\begin{table}[H]
\centering
\caption{Out-of-sample performance on pre-specified walk-forward windows}
\label{tab:mainresults}
\begin{tabularx}{\textwidth}{l *{6}{>{\centering\arraybackslash}X}}
\toprule
Strategy & Ann. return & Ann. vol. & Sharpe & Max drawdown & 95\% CVaR & Turnover \\
\midrule
Equal weight & \todoresult{compute} & \todoresult{compute} & \todoresult{compute} & \todoresult{compute} & \todoresult{compute} & \todoresult{compute} \\
Sample mean--variance & \todoresult{compute} & \todoresult{compute} & \todoresult{compute} & \todoresult{compute} & \todoresult{compute} & \todoresult{compute} \\
DCC-GARCH & \todoresult{compute} & \todoresult{compute} & \todoresult{compute} & \todoresult{compute} & \todoresult{compute} & \todoresult{compute} \\
Rolling vine & \todoresult{compute} & \todoresult{compute} & \todoresult{compute} & \todoresult{compute} & \todoresult{compute} & \todoresult{compute} \\
LSTM-TD3, no vine & \todoresult{compute} & \todoresult{compute} & \todoresult{compute} & \todoresult{compute} & \todoresult{compute} & \todoresult{compute} \\
\methodname & \todoresult{compute} & \todoresult{compute} & \todoresult{compute} & \todoresult{compute} & \todoresult{compute} & \todoresult{compute} \\
\bottomrule
\end{tabularx}
\begin{flushleft}\footnotesize
All returns must be net of the same stated transaction-cost convention. State the annualisation rule, risk-free rate, window aggregation method, number of seeds, and whether figures are pooled across windows or computed per window first.
\end{flushleft}
\end{table}

The primary comparison is economic rather than cosmetic. A higher Sharpe ratio accompanied by implausible turnover, materially worse drawdown, or a gain confined to one window does not establish dominance. Show wealth curves and monthly weights for every window, not only the most favourable one. Since all strategies face the same realised return path within a window, report paired uncertainty: a stationary bootstrap of the return differential or a pre-specified HAC-adjusted test of utility/return differentials is preferable to treating repeated random seeds on one historical path as independent market observations.

\subsection{Ablation and sensitivity evidence}

The ablation table must be generated by retraining and evaluating actual variants under the same locked protocol:
\begin{enumerate}
\item remove $v_t$ while retaining all other features;
\item omit synthetic pre-training but keep the total tuning budget fixed and disclosed;
\item set $\lambda=0$;
\item replace the LSTM encoder with a capacity-matched feed-forward encoder; and
\item compare the corrected monthly scenario generator with any alternative generator.
\end{enumerate}
Each ablation requires its own seeds and walk-forward results. Do not infer component contributions from a single run or from changes in a different benchmark. Sensitivity analysis likewise requires actual retraining or, if computationally infeasible, a clearly defined robustness subset. Hyperparameters must be selected in validation windows, not because they maximise the final holdout result.

\subsection{Interpretation rules}

If the full model wins robustly, the discussion should attribute the result narrowly: the evidence would support the usefulness of its particular dependence-feature representation under this asset universe, frequency, constraints, and cost specification. It would not prove that vine copulas generally improve RL allocation or that synthetic data creates alpha. If performance is neutral or negative, report that finding directly; plausible explanations include weak predictability, feature redundancy, generator misspecification, short effective history, and optimisation instability. Either outcome is scientifically informative when the protocol is clean.
